
\chapter{Предобработка}
Язык Reflex обладает несколькими особенностями из-за которых порождение условий корректности по абстрактному синтаксическому дереву полученному напрямую от программы является не эффективным. Поэтому предполагается несколько этапов предобработки предназначенные для упрощения и оптимезации процесса порождения условий корректности.

\section{Name mangling}
В языке Reflex переменные и константы могут быть объявлены на нескольких уровнях: программы, узла, процесса или кода. При этом работает name shadowing аналогично языку C -- в рамках одного блока возможно единственное обьявление переменной, однако на более глубоком уровне возможно повторное объявление локально перекрывающее более раннее. Также переменные с одинаковыми именами могут быть объявлены в разных процессах и т.д. Все это требует разрешения возможных конфликтов имен.

Сначала определим поля глобального и локального контекста.

Поля глобального контекста:
\begin{itemize}
    \item \textbf{nodeVariables} -- словарь, отображающий имя узла в словарь, отображающий старое имя переменной в новое имя переменной.
    \item \textbf{isDirect} -- словарь, отображающий имя переменной, в булево значение, означающее имеет ли переменная модификатор direct. 
\end{itemize}
Поля локального контекста:
\begin{itemize}
    \item \textbf{path} -- список имен определяющий текущую глубину обработки переменных.
    \item \textbf{variableMap} -- словарь, отображающий старое имя пеерменной в новое.
    \item \textbf{directCount} -- отображающий имя переменной с модификатором direct, в число, означающее количество обращений к нему. 
\end{itemize}

Далее определим функции пребразования имен:
\begin{itemize}
    \item new-name(path, name) $\rightarrow$ \textit{string} -- основная функция конструирования имени, принимает путь path, на котором определена переменная и имя name.
    \item new-indirect-name(name, pos) $\rightarrow$ \textit{string} -- функция конструирующая имя переменной с модификатором indirect, привязанной к битам физического регистра.
    \item new-direct-name(name, num) $\rightarrow$ \textit{string} -- функция конструирующая имя переменной с модификатором direct, привязанной к битам физического регистра.
\end{itemize}

Она принимает один аргумент name - локальное имя переменной:
\begin{lstlisting}[style=PseudoStyle,mathescape=true]
fun newName(path,name){
    new-name = ""
    for p in path{
        new-name += "." + p
    }
    return new-name + "::" + name
}
fun newIndirectName(name,num){
    return name + "_" + num
}
fun newDirectName(name,pos){
    return name + "." + pos
}
\end{lstlisting}

Далее определим функцию mapName(st), где st обозначает обрабатываемую конструкцию программы.
\begin{lstlisting}[style=PseudoStyle,mathescape=true]
//elementAccess - любой вызов переменной
fun mapName(elementAccess) {
    name = elementAccess.name
    accesses = elementAccess.accesses
    newName = ctx.variableMap[name]
    num = ctx.directCount[newName]
    if (env.isDirect[newName]) {
        elementAccess.name = newDirectName(newName, num)
        ctx.directCount[newName] = num + 1
    } else {
        elementAccess.name = newName
    }
    for access in accesses {
        mapName(access)
    }
}

//binaryExpression $\in$ {+,-,*,/,%,>>,<<,==,!=,>=,<=,<,>,&&,||,&,|,^}
fun mapName(binaryExpression) {
    mapName(binaryExpression.left);
    mapName(binaryExpression.right); 
}

//unaryExpression $\in$ {cast,!,-,~}
fun mapName(unaryExpression) {
    expr = unaryExpression.expression
    mapName(expr)
}

//unaryExpression $\in$ {++.,--.,.++,.--}
fun mapName(prefixOrPostfix) {
    access = prefixOrPostfix.access
    mapName(access)
}

//assignment $\in$ {=,+=,-=,*=,/=,%=,<<=,>>=,&=,|=,^=}
fun mapName(assignmentExpression) {
    access = assignmentExpression.left
    accesses = access.accesses
    newName = ctx.variableMap[access.name]
    access.name = newName
    mapName(assignmentExpression.right)
    for access in accesses {
        mapName(access)
    }
}

fun mapName(expressionList) {
    sts = expressionList
    for stmt in sts {
        mapName(stmt)
    }
}

fun mapName(timeoutStatement) {
    savedVariableMap = ctx.variableMap.clone()
    mapName(timeoutStatement.controllingExpression)
    mapName(timeoutStatement.statement)
    ctx.variableMap = savedVariableMap
}

fun mapName(statementList) {
    for stmt in statementList {
        mapName(stmt)
    }
}

fun mapName(statementBlock) {
    ctx.path.clone()
    savedVariableMap = ctx.variableMap.clone()
    ctx.path.pushfront(blockName(uid(statementBlock)))
    mapName(statementBlock.statements)
    ctx.path.pop()
    ctx.variableMap = savedVariableMap
}

fun mapName(waitStatement) {
    mapName(waitStatement.condition)
}

fun mapName(waitOnTimeout) {
    savedVariableMap = ctx.variableMap.clone()
    mapName(waitOnTimeout.condition)
    mapName(waitOnTimeout.controllingExpression)
    mapName(waitOnTimeout.statements)
    ctx.variableMap = savedVariableMap
}

fun mapName(ifThen) {
    mapName(ifThen.condition)
    savedVariableMap = ctx.variableMap.clone()
    mapName(ifThen.then)
    ctx.variableMap = savedVariableMap
}

fun mapName(ifThenElse) {
    mapName(ifThenElse.condition)
    savedVariableMap = ctx.variableMap.clone()
    mapName(ifThenElse.then)
    ctx.variableMap = savedVariableMap
    savedVariableMap = ctx.variableMap.clone()
    mapName(ifThenElse.else)
    ctx.variableMap = savedVariableMap
}

fun mapName(switchS) {
    savedVariableMap = ctx.variableMap.clone()
    ctx.path.pushfront(sblockName(uid(switchS)))
    mapName(switchS.controllingExpression)
    cases = switchS.cases
    for case in cases {
        mapName(case)
    }
    default = switchS.default
    if (default) {
        mapName(default)
    }    
    ctx.path.pop()
    ctx.variableMap = savedVariableMap
}

fun mapName(defaultStatement) {
    mapName(defaultStatement.statements)
}

fun mapName(caseStatement) {
    mapName(caseStatement.statements)
}

fun mapName(expressionS) {
    mapName(expressionS.expression)
}

fun mapName(forS) {
    savedVariableMap = ctx.variableMap.clone()
    mapName(forS.init)
    mapName(forS.condition)
    mapName(forS.update)
    mapName(forS.statement)
    ctx.variableMap = savedVariableMap
}

fun mapName(stateDeclaration) {
    name = stateDeclaration.name
    ctx.path.pushfront(name)
    ctx.path = cons(name.path)
    savedVariableMap = ctx.variableMap.clone()
    sts = stateDeclaration.statements
    for stmt in sts {
        mapName(stmt)
    }
    ctx.path.pop()
    ctx.variableMap = savedVariableMap
    ctx.currentState = nil
}


fun mapName(processDeclaration) {
    savedVariableMap = ctx.variableMap
    node = processDeclaration.node
    nodeVariables = ctx.nodeVariables[node]
    ctx.variableMap = ctx.variableMap $\cup$ nodeVariables
    name = processDeclaration.name
    ctx.path.pushfront(name)
    ctx.currentProcess = name
    states = processDeclaration.states
    for state in states {
        mapName(state)
    }
    ctx.currentProcess = nil
    ctx.path.pop()
    ctx.variableMap = savedVariableMap
}

fun mapName(constantDeclaration) {
    name = constantDeclaration.name
    newName = newName(ctx.path, name)
    constantDeclaration.name = newName
    ctx.variableMap[name] = newName
}

fun mapName(simpleVariableDeclaration) {
    mapName(simpleVariableDeclaration.init)
    name = simpleVariableDeclaration.name
    newName = newName(ctx.path, name)
    simpleVariableDeclaration.name = newName
    ctx.variableMap[name] = newName
}

fun mapName(arrayVariableDeclaration) {
    mapName(arrayVariableDeclaration.init)
    name = arrayVariableDeclaration.name
    newName = newName(ctx.path, name)
    arrayVariableDeclaration.name = newName
    ctx.variableMap[name] = newName
}

fun mapName(physicalVariableDeclaration) {
    name = physicalVariableDeclaration.name
    direct = physicalVariableDeclaration.direct
    newName = newIndirectName(physicalVariableDeclaration.port, physicalVariableDeclaration.index)
    physicalVariableDeclaration.name = newName
    ctx.variableMap[name] = newName
    env.isDirect[newName] = direct
    if (direct) {
        ctx.directCount[newName] = 0
    }
}

fun mapName(enumVariableDeclaration) {
    name = enumVariableDeclaration.name
    newName = newName(ctx.path, name)
    enumVariableDeclaration.name = newName
    ctx.variableMap[name] = newName
}

fun mapName(structureVariableDeclaration) {
    mapName(arrayVariableDeclaration.init)
    name = structureVariableDeclaration.name
    newName = newName(ctx.path, name)
    structureVariableDeclaration.name = newName
    ctx.variableMap[name] = newName
}

fun mapName(nodeDeclaration) {
    variables = nodeDeclaration.variables
    savedVariableMap = ctx.variableMap.clone()
    for varDecl in variables {
        mapName(varDecl)
    }
    ctx.nodeVariables[nodeDeclaration.name] = ctx.variableMap
    ctx.variableMap = savedVariableMap
}

fun mapName(programDeclaration) {
    globalVars = programDeclaration.variables
    nodes = programDeclaration.nodes
    processes = programDeclaration.processes
    for varDecl in globalVars {
        mapName(varDecl)
    }
    for node in nodes {
        mapName(node)
    }
    for process in processes {
        mapName(process)
    }
}
\end{lstlisting}
\section{Проставление cast}
Выражения языка Си во многих случаях производят неявные преобразования типов. В данном блоке описывается алгоритм делающий их явными. Также происводится проставление типов результата каждого выражения.

Сначала определим поля глобального и локального контекста.

Поля глобального контекста:
\begin{itemize}
    \item \textbf{variableType} -- словарь, отображающий имя перемемнной в ее тип.
    \item \textbf{structVariables} -- словарь, отображающий имя структуры в словарь, отображающий имя поля в его тип.
\end{itemize}
Локальный контекст в данном случае полей не имеет.

Определим вспомогательные функции:
\begin{lstlisting}
fun getVariableTypeSup(env, type, fields) {
    if (isInstance(type, "simpleType")) {
        return fields.length == 0 ? type : null;
    }
    if (isInstance(type, "arrayType")) {
        if (typeof fields[0] == "string") {
            return null;
        }
        return getVariableTypeSup(env, type.elementType, fields.slice(1));
    }
    if (isInstance(type, "structureType")) {
        structName = type.name;
        structFields = env.structFields[structName];
        fieldType = structFields[fields[0]];
        if (!fieldType) return null;
        return getVariableTypeSup(env, fieldType, fields.slice(1));
    }
    if (isInstance(type, "enumType")) {
        return fields.length === 0 ? "int32" : null;
    }
    return null; 
}

fun getVariableType(env, name, fields) {
    const baseType = env.variableTypes[name];
    if (!baseType) return null;
    return getVariableTypeSup(env, baseType, fields);
}

function orderIntTypes(left, right) {
    if (left === right) return "=";
    typeList = [
        "bool",
        "int8",  "uint8",
        "int16", "uint16",
        "int32", "uint32",
        "int64", "uint64"
    ];
    idxLeft = typeList.indexOf(left);
    idxRight = typeList.indexOf(right);
    if (idxLeft == idxRight) {
        return "=";
    }
    return idxLeft < idxRight ? "<" : ">";
}

function defType(op, left, right) {
    if (isInstance(left, "undefinedType")) {
        if (isInstance(right, "undefinedType")) {
            if (isInstance(right, "undefinedFloatType") ||
                isInstance(left,  "undefinedFloatType")) {
                return "undefinedFloatType";
            }
            return "undefinedIntType";
        }
        return right;
    }
    if (isInstance(right, "undefinedType")) {
        return left;
    }
    if (isInstance(left, "floatType") || isInstance(right, "floatType")) {
        if (left == "double" || right == "double") {
            return "double";
        }
        return "float";
    }
    switch (op) {
        case "+": case "-": case "*": case "/": case "%":
        case ">>": case "<<":
        case "==": case "!=": case ">=": case "<=": case ">": case "<":
            intl = orderIntTypes(left, "int32");
            intr = orderIntTypes(right, "int32");
            comp = orderIntTypes(left, right);
            if (intl == "<" && intr == "<") return "int32";
            return comp == "<" ? right : left;
        case "&&": case "||":
            return "bool";
        case "&": case "|": case "^":
            return orderIntTypes(left, right) == "<" ? right : left;
        case "!.":
            return "bool";
        case "-." : case "~." :
            const intr = orderIntTypes(right, "int32");
            return intr == ">" ? right : "int32";
        case "=": case "+=": case "-=": case "*=": case "/=":
        case "<<=": case ">>=":
            intl2 = orderIntTypes(left, "int32");
            intr2 = orderIntTypes(right, "int32");
            comp2 = orderIntTypes(left, right);
            if (intl2 == "<" && intr2 == "<") return "int32";
            return comp2 == "<" ? right : left;
        case "&=": case "|=": case "^=":
            return orderIntTypes(left, right) == "<" ? right : left;
        case "++." : case ".++" : case "--." : case ".--" :
            return right;
        default:
            return left;
    }
}
\end{lstlisting}

Далее определим функцию specifyType(st), где st обозначает обрабатываемую конструкцию программы.

\begin{lstlisting}
fun specifyType(boolConstant){
    return "bool"
}
fun specifyType(integerConstant){
    return "undefinedIntType"
}
fun specifyType(naturalConstant){
    return "undefinedIntType"
}
fun specifyType(floatConstant){
    return "undefinedFloatType"
}
fun specifyType(timeConstant){
    return "undefinedIntType"
}

fun specifyType(elementAccess) {
    name = elementAccess.variable
    accesses = elementAccess.accesses
    if (!accesses) {
        ty = getVariableType(env, name, [])
        elementAccess.restype = ty
        return ty
    }
    ty = getVariableType(env, name, [])
    path = []
    for access in accesses {
        ty = getVariableType(env, name, reverse(path))
        
        if (access.type == "field name") {
            path.pushfront(access)
        } 
        else if (access.type == "expression") {
            specifyType(access)
            path.pushfront(0)
        }
    }
    ty = getVariableType(env, name, reverse(path))
    elementAccess.restype = ty
    return ty
}
//binaryBoolExpression $\in$ {==,!=,>=,<=,>,<,&&,||}
fun specifyType(binaryBoolExpression) {
    left = specifyType(binaryBoolExpression.left)
    right = specifyType(binaryBoolExpression.right)
    newType = defType(otype(binaryBoolExpression), left, right)
    if (newType != left) {
        if (!isInstance(left, "undefinedType") || isInstance(newType, "undefinedType")) {
            binaryBoolExpression.left = createCast(
                type: newType,
                expression: binaryBoolExpression.left,
                pretype: left
            )
        }
    }
    if (newType != right) {
        if (!isInstance(right, "undefinedType") || isInstance(newType, "undefinedType")) {
            binaryBoolExpression.right = createCast(
                type: newType,
                expression: binaryBoolExpression.right,
                pretype: right
            )
        }
    }
    binaryBoolExpression.restype = "bool"
    return "bool"
}

//binaryBoolExpression $\in$ {+,-,*,/,>>,<<,&,|,^}
fun specifyType(binaryExpression) {
    left  = specifyType(binaryExpression.left)
    right = specifyType(binaryExpression.right)
    leftArg  = binaryExpression.left
    rightArg = binaryExpression.right
    newType = defType(otype(binaryExpression), left, right)
    if (newType != left) {
        if (!isInstance(left, "undefinedType") || isInstance(newType, "undefinedType")) {
            binaryExpression.left = createCast(
                type: newType,
                expression: leftArg,
                pretype: left
            )
        }
    }
    if (newType != right) {
        if (!isInstance(right, "undefinedType") || isInstance(newType, "undefinedType")) {
            binaryExpression.right = createCast(
                type: newType,
                expression: rightArg,
                pretype: right
            )
        }
    }
    binaryExpression.restype = newType
    return newType
}

fun specifyType(castExpression) {
    specifyType(castExpression.expression)
    castExpression.restype = castExpression.type
    return castExpression.type
}
//unaryExpression $\in$ {!,-,~}
fun specifyType(unaryExpression) {
    operandType = specifyType(unaryExpression.right)
    newType = defType(otype(unaryExpression), null, operandType)
    if (newType != operandType) {
        if (!isInstance(operandType, "undefinedType") || 
            isInstance(newType, "undefinedType")) {
            unaryExpression.right = createCast(
                type: newType,
                expression: unaryExpression.right,
                pretype: operandType
            )
        }
    }
    unaryExpression.restype = newType
    return newType
}

fun specifyType(assignmentExpression) {
    rightType = specifyType(assignmentExpression.right)
    leftType = specifyType(assignmentExpression.left)
    if (leftType != rightType && !isInstance(rightType, "undefinedType")) {
        assignmentExpression.right = createCast(
            type: leftType,
            expression: assignmentExpression.right,
            pretype: rightType
        )
    }
    assignmentExpression.restype = leftType
    return leftType
}

fun specifyType(prefixOrPostfix) {
    val = specifyType(prefixOrPostfix.access)
    prefixOrPostfix.restype = val
    return val
}

fun specifyType(processStateChecking) {
    processStateChecking.restype = "bool"
    return "bool"
}

fun specifyType(statementList) {
    for stmt in statementList {
        specifyType(stmt)
    }
}

fun specifyType(timeoutStatement) {
    specifyType(timeoutStatement.statements)
}

fun specifyType(waitStatement) {
    specifyType(waitStatement.condition)
}

fun specifyType(waitOnTimeout) {
    specifyType(waitOnTimeout.condition)
    specifyType(waitOnTimeout.statements)
    specifyType(waitOnTimeout.controllingExpression)
}

fun specifyType(ifThen) {
    condType = specifyType(ifThen.condition)
    
    if (!isInstance(condType, "booleanType")) {
        ifThen.condition = createCast(
            type: "bool",
            expression: ifThen.condition,
            pretype: condType
        )
    } 
    specifyType(ifThen.then)
}

fun specifyType(ifThenElse) {
    condType = specifyType(ifThenElse.condition)
    if (!isInstance(condType, "booleanType")) {
        ifThenElse.condition = createCast(
            type: "bool",
            expression: ifThenElse.condition,
            pretype: condType
        )
    }
    specifyType(ifThenElse.then)
    specifyType(ifThenElse.else)
}

fun specifyType(switchS) {
    specifyType(switchS.controllingExpression)
    for caseStmt in switchS.cases {
        specifyType(caseStmt)
    }
    if (switchS.default) {
        specifyType(switchS.default)
    }
}

fun specifyType(caseStatement) {
    specifyType(caseStatement.statements)
}

fun specifyType(defaultStatement) {
    specifyType(defaultStatement.statements)
}

fun specifyType(statementBlock) {
    specifyType(statementBlock.statements)
}

fun specifyType(expressionS) {
    specifyType(expressionS.expression)
}

fun specifyType(forS) {
    specifyType(forS.init)
    specifyType(forS.condition)
    specifyType(forS.update)
    specifyType(forS.statement)
}

fun specifyType(constantDeclaration) {
    env.variableTypes[constantDeclaration.name] = constantDeclaration.type
}

fun specifyType(simpleVariableDeclaration) {
    env.variableTypes[simpleVariableDeclaration.name] = simpleVariableDeclaration.type
}

fun specifyType(arrayVariableDeclaration) {
    env.variableTypes[arrayVariableDeclaration.name] = arrayVariableDeclaration.type
}

fun specifyType(physicalVariableDeclaration) {
    env.variableTypes[physicalVariableDeclaration.name] = physicalVariableDeclaration.type
}

fun specifyType(enumVariableDeclaration) {
    env.variableTypes[enumVariableDeclaration.name] = enumVariableDeclaration.type
}

fun specifyType(structureVariableDeclaration) {
    env.variableTypes[structureVariableDeclaration.name] = structureVariableDeclaration.type
}

fun specifyType(structureDeclaration) {
    fieldTypes = {}
    for field in structureDeclaration.fields {
        fieldTypes[fieldName] = fieldType
    }
    env.structFields[structureDeclaration.name] = fieldTypes
}

fun specifyType(nodeDeclaration) {
    for varDecl in nodeDeclaration.variables{
        specifyType(varDecl)
    }
}

fun specifyType(stateDeclaration) {
    specifyType(stateDeclaration.statements)
}

fun specifyType(processDeclaration) {
    for state in processDeclaration.states {
        specifyType(state)
    }
}

fun specifyType(programDeclaration) {
    for node in programDeclaration.nodes {
        specifyType(node)
    }
    for process in programDeclaration.processes {
        specifyType(process)
    }
}
\end{lstlisting}

\section{Нормализация}
Нормализация заключается в преобразовании некоторых конструкций в программе для упрощения анализа и обработки программы.

Сначала определим поля глобального и локального контекста.

Поля глобального контекста:
\begin{itemize}
    \item \textbf{processStatesNames} -- словарь, отображающий имя процесса  в список имен состояний
\end{itemize}

Локальный контекст имеет следующие поля:
    \begin{itemize}
        \item \textbf{currentProcess} -- имя текущего процесса.
        \item \textbf{currentState} -- имя текущего состояния.
        \item \textbf{nextState} -- имя следующего состояния.
        \item \textbf{lightСounter} -- счетчик оператора числа лекговесных состояний в процессе.
    \end{itemize}
    
Определим вспомогательные функции:
\begin{lstlisting}
fun nextProcessState() {
    statesList = env.processStateNames[ctx.currentProccess]
    if (statesList == null) {
        return nil
    }
    idx = statesList.indexOf(ctx.currentState)
    if (idx == -1 || idx == statesList.length - 1) {
        return nil
    }
    return statesList[idx + 1]
}

fun firstState(process) {
    states = env.processStatesNames[process]
    if (states == null || states.length == 0) {
        return nil
    }
    return states[0]
}

fun createLightState(stateName, stateType, num) {
    return stateName + "_" + stateType + "_" + num
}

fun breakIntoStates(statements) {
    result = []
    currentState = []
    for stmt in statements {
        if (isInstance(stmt, "slice") ||
            isInstance(stmt, "wait") ||
            isInstance(stmt, "waitOnTimeout")) {
            if (currentState.length > 0) {
                result.pushfront(reverse(currentState))
                currentState = []
            }
            currentState = [stmt]
        } else {
            currentState.pushfront(stmt)
        }
    }
    if (currentState.length > 0) {
        result.pushfront(reverse(currentState))
    }
    return reverse(result)
}
\end{lstlisting}

Далее определим функцию normalizeProgram(st), где st обозначает обрабатываемую конструкцию программы. В большенстве случаев, она ничего не возвращает, исключение - обработка stateDeclaration, после которой возвращается список новых состояний.
\begin{lstlisting}
fun normalizeProgram(setState) {
    setState.state = ctx.nextState
}

fun normalizeProgram(statementList) {
    for stmt in statementList {
        normalizeProgram(stmt)
    }
}

fun normalizeProgram(switchS) {
    cases = switchS.cases
    default = switchS.sefault
    for case in cases {
        normalizeProgram(case)
    }
    normalizeProgram(default)
    excases = []
    for case in cases {
        for excase in excases {
            excase.statements = excase.statements + case.statements.clone() 
        }
        if (!case.break){
            excases.pushfront(case)
        } else{
            for excase in excases {
                excase.break = true
            }
            excases = []
        }
        
    }
    if (default){
        for excase in exCases {
            excase.statements = excase.statements + default.statements 
        }
        for excase in excases {
                excase.break = true
            }
    }
}

fun normalizeProgram(caseStatement) {
    normalizeProgram(caseStatement.statements)
}

fun normalizeProgram(defaultStatement) {
    normalizeProgram(defaultStatement.statements)
}

fun normalizeProgram(blockStatement) {
    normalizeProgram(blockStatement.statements)
}

fun normalizeProgram(ifThen) {
    normalizeProgram(ifThen.then)
}

fun normalizeProgram(ifThenElse) {
    normalizeProgram(ifThenElse.then)
    normalizeProgram(ifThenElse.else)
}

fun normalizeProgram(forS) {
    if (forS.init)      normalizeProgram(forS.init)
    if (forS.condition) normalizeProgram(forS.condition)
    if (forS.update)    normalizeProgram(forS.update)
    if (forS.statement) normalizeProgram(forS.statement)
}

fun normalizeState(name,sts){
    if (sts.empty()){
        return {body:[],newStates:[],mustTransfer:false}
    }
    head = sts.pop()
    normalizeProgram(head)
    tail = sts 
    tailResult = normalizeState(name,tail)

    tailBody = tailResult.body
    tailStates = tailResult.newStates
    tailTransfer = tailResult.mustTransfer

    if (isInstance(head,"slice")){
        counter = ctx.lightСounter
        ctx.lightСounter++
        newName = name + "_light" + counter

        newState = createState({name:newName, statements:tailBody})

        return {body: createSetState(newName), 
            newStates: tailStates $\cup$ [newState],
            mustTransfer = true}
    } else
    if (isInstance(head,"wait")){
        counter = ctx.lightСounter
        ctx.lightСounter+=2
        newName1 = name + "_light" + counter
        newName1 = name + "_light" + (counter + 1)
        if (tailTransfer){
            thenBody = tailBody
            ifStatement = createIfThenElse({
                condition: head.condition, 
                then: thenBody, 
                else: [createSetState(newName1)]})
            state1 = createState({
                name:newName1, 
                statements: [createifThen({condition: head.condition, then: thenBody})]})
            return {body: [ifStatement],
                newStates: tailStates $\cup$ [state1],
                mustTransfer:True}
        }else{
            thenBody = tailBody.pushback(createSetState)
            ifStatement = createIfThenElse({
                condition: head.condition, 
                then: thenBody, 
                else: [createSetState(newName1)]})
            state1 = createState({
                name:newName1, 
                statements: [createifThen({condition: head.condition, then: thenBody})]})
            state2 = createState({
                name:newName2,
                statements: tailBody
            })
            return {body: [ifStatement],
                newStates: tailStates $\cup$ [state1,state2],
                mustTransfer:True} 
        }
    } else 
    if (isInstance(head,"waitOnTimeout")){
        counter = ctx.lightСounter
        ctx.lightСounter+=2
        newName1 = name + "_light" + counter
        newName1 = name + "_light" + (counter + 1)
        if (tailTransfer){
            thenBody = tailBody
            ifStatement = createIfThenElse({
                condition: head.condition, 
                then: thenBody, 
                else: [createSetState(newName1)]})
            state1 = createState({
                name:newName1, 
                statements: [createifThen({condition: head.condition, then: thenBody}),
                    createTimeoutStatement({
                        controllingExpression: head.controllingExpression,
                        statements: head.statements
                    })]})
            return {body: [ifStatement],
                newStates: tailStates $\cup$ [state1],
                mustTransfer:True}
        }else{
            thenBody = tailBody.pushback(createSetState)
            ifStatement = createIfThenElse({
                condition: head.condition, 
                then: thenBody, 
                else: [createSetState(newName1)]})
            state1 = createState({
                name:newName1, 
                statements: [createifThen({condition: head.condition, then: thenBody}),
                    createTimeoutStatement({
                        controllingExpression: head.controllingExpression,
                        statements: head.statements
                    })]})
            state2 = createState({
                name:newName2,
                statements: tailBody
            })
            return {body: [ifStatement],
                newStates: tailStates $\cup$ [state1,state2],
                mustTransfer:True} 
        }
    } else{
        return {body: tailBody.pushfront(head),
            newStates: tailStates,
            mustTransfer: tailTransfer}
    }
}

fun normalizeProgram(stateDeclaration){
    name = stateDeclaration.name
    sts = stateDeclaration.statements
    res = normalizeState(name, sts)
    stateDeclaration.sts = res.body
    res.newStates.pushfront(stateDeclaration)
    return res.newStates
}

fun normalizeProgram(processDeclaration) {
    name = processDeclaration.name
    states = processDeclaration.states

    ctx.currentProcess = name
    newStates = []
    for state in states {
        res = normalizeProgram(state)
        newstates.pushback(res)
    }
    processDeclaration.states = newStates
}


fun normalizeProgram(programDeclaration) {
    processes = programDeclaration.processes
    for proc in processes {
        normalizeProgram(proc)
    }
}
\end{lstlisting}

